(a) The only nonempty open subsets are $X$ and its generic point $\eta$. Thus an $\mathcal O_X$-module is specified by an $R$-module $M$, a $K$-vector space $L$, and an $R$-linear restriction $M\to L$. By localization, such restrictions correspond to $K$-linear maps $\rho:M\otimes_RK\to L$. There are no further sheaf conditions, since every open covering of a nonempty open set contains that open set.

(b) The sheaf $\widetilde M$ has sections $M$ on $X$ and $M\otimes_RK$ on $\eta$, with the localization restriction. Since every quasi-coherent sheaf on $X$ is associated to its module of global sections, the given sheaf is quasi-coherent exactly when $\rho$ is an isomorphism.
